\documentclass[11pt]{article}

% ---------- FangMath house style ----------
\usepackage[T1]{fontenc}
\usepackage[utf8]{inputenc}
\usepackage{helvet}
\renewcommand{\familydefault}{\sfdefault}
\usepackage[margin=1.9cm,top=2.4cm,bottom=2.4cm]{geometry}
\usepackage{xcolor}
\usepackage[most]{tcolorbox}
\usepackage{fancyhdr}
\usepackage{titlesec}
\usepackage{enumitem}
\usepackage{amsmath}
\usepackage{amssymb}
\usepackage{array}
\usepackage{needspace}
\usepackage{multicol}
\usepackage{tikz}
\usetikzlibrary{arrows.meta,calc}
\usepackage{hyperref}
\hypersetup{hidelinks}

\definecolor{fmnavy}{RGB}{0,32,96}
\definecolor{fmgray}{RGB}{110,110,110}
\definecolor{keybg}{RGB}{224,242,241}
\definecolor{keyframe}{RGB}{0,105,92}
\definecolor{workbg}{RGB}{250,250,250}
\definecolor{workframe}{RGB}{190,190,190}
\definecolor{exbg}{RGB}{232,245,233}
\definecolor{exframe}{RGB}{46,125,50}
\definecolor{mistbg}{RGB}{253,235,235}
\definecolor{mistframe}{RGB}{183,28,28}
\definecolor{ytbg}{RGB}{255,248,225}
\definecolor{ytframe}{RGB}{191,144,0}

\titleformat{\section}
  {\color{fmnavy}\Large\bfseries}{\thesection.}{0.6em}{}
\titlespacing*{\section}{0pt}{16pt}{4pt}

\newcommand{\sectionepigraph}[1]{{\color{fmgray}\itshape\small #1\par}\vspace{6pt}}

\pagestyle{fancy}
\fancyhf{}
\renewcommand{\headrulewidth}{0.6pt}
\renewcommand{\footrulewidth}{0pt}
\fancyhead[L]{\small\color{fmnavy}\bfseries FangMath}
\fancyhead[R]{\small\color{fmgray} Real Numbers \& Algebraic Expressions}
\fancyfoot[C]{\small\color{fmgray} \textcopyright\ 2026 FangMath. All rights reserved.}
\fancyfoot[R]{\small\color{fmgray}\thepage}

\setlength{\parindent}{0pt}
\setlength{\parskip}{4pt}

\newcommand{\infoblank}[1]{\underline{\hspace{#1}}}

\newtcolorbox{keyidea}[1][Key Idea]{colback=keybg,colframe=keyframe,
  fonttitle=\bfseries,coltitle=white,colbacktitle=keyframe,title=#1,
  arc=1mm,boxrule=0.6pt}
\newtcolorbox{workedex}[1][Watch Me]{colback=exbg,colframe=exframe,
  fonttitle=\bfseries,coltitle=white,colbacktitle=exframe,title=#1,
  arc=1mm,boxrule=0.6pt}
\newtcolorbox{mistakebox}[1][Common Mistake to Avoid]{colback=mistbg,colframe=mistframe,
  fonttitle=\bfseries,coltitle=white,colbacktitle=mistframe,title=#1,
  arc=1mm,boxrule=0.6pt}

% ---------- Ruled student workspace (responsive width) ----------
\newcommand{\ruleline}{\vspace{0.55cm}\noindent\textcolor{workframe}{\hrulefill}\par}
\newcommand{\drawrulelines}[1]{%
  \count255=#1\relax
  \loop
    \ruleline
    \advance\count255 by -1
  \ifnum\count255>0
  \repeat
}
\newcommand{\workspace}[1]{%
  \begin{tcolorbox}[colback=workbg, colframe=workframe, boxrule=0.4pt, arc=1mm,
                     left=6pt, right=6pt, top=2pt, bottom=4pt]
    {\footnotesize\color{fmgray} Show your work:}
    \drawrulelines{#1}
  \end{tcolorbox}
}

\newcounter{prac}
\newcommand{\practice}[1]{%
  \needspace{4\baselineskip}
  \refstepcounter{prac}
  \par\vspace{4pt}{\bfseries\color{fmnavy} \theprac.} #1\par
}

\begin{document}

% =====================================================================
% TITLE BLOCK
% =====================================================================
\begin{center}
  {\Huge\color{fmnavy}\bfseries Real Numbers \& Algebraic Expressions}\\[6pt]
  {\large\color{fmgray}\itshape A Chapter 1 practice worksheet --- classifying, comparing, and\break
  working with real numbers, plus writing and evaluating expressions.}
\end{center}

\vspace{12pt}

\begin{tcolorbox}[colback=workbg, colframe=workframe, boxrule=0.4pt, arc=1mm,
                   left=14pt, right=14pt, top=10pt, bottom=10pt]
  \renewcommand{\arraystretch}{2.2}
  \begin{tabular}{@{}l l@{}}
    {\bfseries\color{fmnavy} Name:}~\infoblank{5.4cm} &
    {\bfseries\color{fmnavy} Grade:}~\infoblank{2.6cm} \\
    {\bfseries\color{fmnavy} Date:}~\infoblank{3.4cm} &
    {\bfseries\color{fmnavy} Weekday:}~\infoblank{3.4cm} \\
    {\bfseries\color{fmnavy} Time:}~\infoblank{3.4cm} &
  \end{tabular}
\end{tcolorbox}

\vspace{6pt}

% =====================================================================
\section{A. Classifying Real Numbers}
\sectionepigraph{Every real number belongs to at least one of these families. Some numbers belong to several.}

\begin{keyidea}[The Real Number Family]
\textbf{Natural numbers} ($1,2,3,\dots$) are the counting numbers.\\
\textbf{Whole numbers} ($0,1,2,3,\dots$) are the natural numbers plus $0$.\\
\textbf{Integers} ($\dots,-2,-1,0,1,2,\dots$) are the whole numbers plus their opposites.\\
\textbf{Rational numbers} can be written as a fraction $\dfrac{a}{b}$ of two integers
($b\neq0$) --- this includes all integers, all terminating decimals, and all
repeating decimals.\\
\textbf{Irrational numbers} \emph{cannot} be written as a fraction of two integers ---
their decimals go on forever with no repeating pattern (like $\sqrt2$ or $\pi$).
Every real number is rational \emph{or} irrational, never both.
\end{keyidea}

\begin{mistakebox}
A square root is not automatically irrational! $\sqrt{36}=6$ is a whole number,
because $36$ is a \emph{perfect square}. Always simplify a square root first ---
\emph{then} decide which family it belongs to.
\end{mistakebox}

\textbf{Name every subset (Natural, Whole, Integer, Rational, Irrational) that each number belongs to.}

\begin{multicols}{2}
\practice{$9$}
\practice{$-14$}
\practice{$\dfrac{3}{8}$}
\practice{$\sqrt{36}$}
\practice{$\sqrt{28}$}
\practice{$-4.75$}
\practice{$0$}
\practice{$\sqrt5$}
\end{multicols}

% =====================================================================
\section{B. Simplifying Square Roots}
\sectionepigraph{If the number under the root sign is a perfect square (or a fraction of two perfect squares), the square root simplifies to an exact number.}

\begin{workedex}
Simplify $\sqrt{\dfrac{49}{100}}$.
\[
\sqrt{\dfrac{49}{100}} = \dfrac{\sqrt{49}}{\sqrt{100}} = \dfrac{7}{10},
\qquad\text{since } 7^2=49 \text{ and } 10^2=100.
\]
\end{workedex}

\practice{$\sqrt{64}$}
\practice{$\sqrt{121}$}
\practice{$\sqrt{\dfrac{9}{16}}$}
\practice{$\sqrt{\dfrac{25}{36}}$}

% =====================================================================
\section{C. Estimating Square Roots}
\sectionepigraph{When a number is \emph{not} a perfect square, its square root sits between two whole numbers. Find which two, then decide which one it's closer to.}

\begin{workedex}
Estimate $\sqrt{52}$ to the nearest integer.
\[
6^2=36 \qquad 7^2=49 \qquad 8^2=64
\]
Since $49<52<64$, we know $7<\sqrt{52}<8$. Because $52$ is much closer to
$49$ than to $64$, $\sqrt{52}\approx 7$.
\end{workedex}

\textbf{Estimate each square root. Round to the nearest integer.}

\practice{$\sqrt{40}$}
\workspace{2}
\practice{$\sqrt{55}$}
\workspace{2}
\practice{$\sqrt{95}$}
\workspace{2}
\practice{$\sqrt{150}$}
\workspace{2}

% =====================================================================
\section{D. Square Roots in Context}
\sectionepigraph{If you know the \emph{area} of a square, its side length is the square root of the area --- since $A = s^2$.}

\practice{A framed painting is a square with an area of $50$ in$^2$. About how long is one side, to the nearest inch?}
\workspace{3}

\practice{A garden plot is a square with an area of $95$ m$^2$. About how long is one side, to the nearest meter?}
\workspace{3}

\practice{A stage floor is a square with an area of $200$ ft$^2$. About how long is one side, to the nearest foot?}
\workspace{3}

% =====================================================================
\section{E. Comparing Real Numbers}
\sectionepigraph{To compare two real numbers, rewrite both as decimals --- then line them up and compare.}

\begin{workedex}
Compare $3\dfrac{1}{4}$ and $\sqrt{11}$ using $<,\ >,$ or $=$.
\[
3\tfrac14 = 3.25, \qquad \sqrt{11}\approx3.3166\dots
\]
Since $3.25<3.3166\dots$, \quad $3\dfrac14 < \sqrt{11}$.
\end{workedex}

\textbf{Compare the numbers in each pair using $<,\ >,$ or $=$.}

\practice{$4\dfrac25 \quad \bigcirc \quad \sqrt{19}$}
\workspace{2}
\practice{$-2.1 \quad \bigcirc \quad -\dfrac74$}
\workspace{2}
\practice{$\dfrac53 \quad \bigcirc \quad \sqrt3$}
\workspace{2}
\practice{$\sqrt{70} \quad \bigcirc \quad 8.3$}
\workspace{2}

% =====================================================================
\section{F. Ordering Real Numbers}
\sectionepigraph{Convert every number to a decimal first --- then it's just like ordering decimals.}

\practice{Order these numbers from \textbf{least to greatest}:
\[
3, \qquad -1, \qquad \sqrt{10}, \qquad -\dfrac52, \qquad 2.6
\]}
\workspace{4}

% =====================================================================
\section{G. Properties of Real Numbers}
\sectionepigraph{These two properties describe ways you can rearrange a sum or a product without changing the answer.}

\begin{keyidea}[Commutative \& Associative Properties]
\textbf{Commutative Property} --- you can swap the \emph{order} of the numbers:
\[
a+b=b+a \qquad\qquad a\times b = b\times a
\]
\textbf{Associative Property} --- you can move the \emph{grouping} parentheses, as
long as the order of the numbers stays the same:
\[
(a+b)+c = a+(b+c) \qquad\qquad (a\times b)\times c = a\times(b\times c)
\]
A quick way to tell them apart: if the \emph{numbers switch places}, it's
Commutative. If the numbers \emph{stay in the same order} but the parentheses move,
it's Associative.
\end{keyidea}

\textbf{Name the property shown in each equation.}

\practice{$(4+9)+6 = 4+(9+6)$}
\practice{$7\times12 = 12\times7$}
\practice{$5\times(3\times8) = (5\times3)\times8$}
\practice{$n+15 = 15+n$}

% =====================================================================
\section{H. Writing Algebraic Expressions}
\sectionepigraph{Turn a word phrase into symbols by translating one piece at a time.}

\begin{workedex}
Write an algebraic expression for: \emph{4 more than three times a number $t$.}
\[
\text{``three times a number } t\text{''} = 3t \qquad\qquad
\text{``4 more than that''} = 3t+4.
\]
\end{workedex}

\practice{A number $k$ increased by $9$.}
\practice{$5$ less than twice a number $m$.}
\practice{The quotient of a number $p$ and $7$.}
\practice{Renting a kayak costs an \$8 flat fee plus \$3 for every hour $h$ you keep it.}

\practice{The table shows the total cost of catering an event based on the number
of guests. Write a rule for the total cost \textbf{in words}, and then as an
algebraic expression using $n$ for the number of guests.

\begin{center}
\renewcommand{\arraystretch}{1.4}
\begin{tabular}{|c|c|}
\hline
\textbf{Number of Guests} & \textbf{Total Cost} \\ \hline
$10$ & $(15\cdot10)+40$ \\ \hline
$20$ & $(15\cdot20)+40$ \\ \hline
$30$ & $(15\cdot30)+40$ \\ \hline
\end{tabular}
\end{center}
}
\workspace{3}

% =====================================================================
\section{I. Order of Operations}
\sectionepigraph{Work inside parentheses first, then exponents, then multiply/divide left to right, then add/subtract left to right --- \emph{PEMDAS}.}

\begin{workedex}
Simplify $30\div(1+2^2)$.
\[
30\div(1+2^2) = 30\div(1+4) = 30\div5 = 6.
\]
\end{workedex}

\practice{$36\div(2+4^2)$}
\workspace{2}
\practice{$5^2-3\times(8-2)$}
\workspace{2}
\practice{$(9-4)^2+3\times2$}
\workspace{2}
\practice{$48\div(2^3-2)+5$}
\workspace{2}

% =====================================================================
\section{J. Evaluating Algebraic Expressions}
\sectionepigraph{Substitute the given value for each variable --- then follow order of operations.}

\begin{workedex}
Evaluate $2x^2\div y$ when $x=3$ and $y=6$.
\[
2(3)^2\div6 = 2(9)\div6 = 18\div6 = 3.
\]
\end{workedex}

\practice{Evaluate $4x\cdot3\div y$ when $x=2$ and $y=6$.}
\workspace{3}

\practice{Evaluate $(3a)^2\div(b-1)$ when $a=2$ and $b=5$.}
\workspace{3}

\practice{Evaluate $2(x+y)-z$ when $x=4$, $y=3$, and $z=5$.}
\workspace{3}

% =====================================================================
\section{K. Apply: Word Problems}
\sectionepigraph{Read carefully, decide which idea from this worksheet applies, and show your reasoning.}

\practice{A contractor is installing tile in a square room. Tile and labor
cost \$8 per square foot, and the total job costs \$960. About how long is a
side of the room, to the nearest foot?}
\workspace{4}

\practice{If you lean a ladder against a wall, a safe length for the ladder is
$\sqrt{x^2+(3x)^2}$ feet, where $x$ is how far the base of the ladder sits from
the wall. Find the ladder length, to the nearest tenth of a foot, when the base
is $6$ feet from the wall.}
\workspace{4}

\practice{A rafting company charges \$48 per adult, \$27 per child, and \$6 per
hour for equipment rental. Write an expression for the total cost of a trip
for $2$ adults and $3$ children who rent equipment for $h$ hours. Then find the
total cost if the trip lasts $4$ hours.}
\workspace{5}

% =====================================================================
\newpage
\section*{\color{fmnavy} Detailed Answer Key}
\sectionepigraph{Full solutions and reasoning, in order.}

\textbf{Section A --- Classifying Real Numbers.}\\
1. $9$: Natural, Whole, Integer, Rational.\\
2. $-14$: Integer, Rational (not Natural or Whole --- it's negative).\\
3. $\dfrac38$: Rational only.\\
4. $\sqrt{36}=6$: Natural, Whole, Integer, Rational (it's a perfect square!).\\
5. $\sqrt{28}\approx5.2915\dots$: Irrational ($28$ is not a perfect square, so the decimal never ends or repeats).\\
6. $-4.75$: Rational (a terminating decimal, and $-4.75=-\dfrac{19}{4}$).\\
7. $0$: Whole, Integer, Rational (not Natural --- the natural numbers start at $1$).\\
8. $\sqrt5\approx2.2360\dots$: Irrational ($5$ is not a perfect square).

\textbf{Section B --- Simplifying Square Roots.}\\
9. $\sqrt{64}=8$, since $8^2=64$.\\
10. $\sqrt{121}=11$, since $11^2=121$.\\
11. $\sqrt{\dfrac{9}{16}}=\dfrac{\sqrt9}{\sqrt{16}}=\dfrac{3}{4}$.\\
12. $\sqrt{\dfrac{25}{36}}=\dfrac{\sqrt{25}}{\sqrt{36}}=\dfrac{5}{6}$.

\textbf{Section C --- Estimating Square Roots.}\\
13. $\sqrt{40}$: since $36<40<49$, $6<\sqrt{40}<7$; $40$ is close to $36$, so $\sqrt{40}\approx6$ ($\sqrt{40}\approx6.32$).\\
14. $\sqrt{55}$: since $49<55<64$, $7<\sqrt{55}<8$; $55$ is closer to $49$, so $\sqrt{55}\approx7$ ($\sqrt{55}\approx7.42$).\\
15. $\sqrt{95}$: since $81<95<100$, $9<\sqrt{95}<10$; $95$ is much closer to $100$, so $\sqrt{95}\approx10$ ($\sqrt{95}\approx9.75$).\\
16. $\sqrt{150}$: since $144<150<169$, $12<\sqrt{150}<13$; $150$ is much closer to $144$, so $\sqrt{150}\approx12$ ($\sqrt{150}\approx12.25$).

\textbf{Section D --- Square Roots in Context.}\\
17. $A=s^2=50 \Rightarrow s=\sqrt{50}\approx7.07$, so the side is about $\mathbf{7}$ inches.\\
18. $A=s^2=95 \Rightarrow s=\sqrt{95}\approx9.75$, so the side is about $\mathbf{10}$ meters.\\
19. $A=s^2=200 \Rightarrow s=\sqrt{200}\approx14.14$, so the side is about $\mathbf{14}$ feet.

\textbf{Section E --- Comparing Real Numbers.}\\
20. $4\dfrac25=4.4$ and $\sqrt{19}\approx4.359$, so $4\dfrac25 \mathbf{>} \sqrt{19}$.\\
21. $-2.1$ and $-\dfrac74=-1.75$; since $-2.1$ is farther left on the number line, $-2.1 \mathbf{<} -\dfrac74$.\\
22. $\dfrac53\approx1.667$ and $\sqrt3\approx1.732$, so $\dfrac53 \mathbf{<} \sqrt3$.\\
23. $\sqrt{70}\approx8.367$ and $8.3$, so $\sqrt{70} \mathbf{>} 8.3$.

\textbf{Section F --- Ordering Real Numbers.}\\
24. Converting: $3=3.000$, $-1=-1.000$, $\sqrt{10}\approx3.162$, $-\dfrac52=-2.5$, $2.6=2.600$.\\
From least to greatest: $\mathbf{-\dfrac52,\ -1,\ 2.6,\ 3,\ \sqrt{10}}$.

\textbf{Section G --- Properties of Real Numbers.}\\
25. $(4+9)+6=4+(9+6)$: the order of the numbers doesn't change, only the grouping --- \textbf{Associative Property of Addition}.\\
26. $7\times12=12\times7$: the two numbers switch places --- \textbf{Commutative Property of Multiplication}.\\
27. $5\times(3\times8)=(5\times3)\times8$: the order stays the same, only the grouping moves --- \textbf{Associative Property of Multiplication}.\\
28. $n+15=15+n$: the two terms switch places --- \textbf{Commutative Property of Addition}.

\textbf{Section H --- Writing Algebraic Expressions.}\\
29. $k+9$.\\
30. $2m-5$.\\
31. $\dfrac{p}{7}$ (or $p\div7$).\\
32. $8+3h$.\\
33. Each cost is \$15 times the number of guests, plus a \$40 fee. In words:
``the total cost is \$15 per guest plus a \$40 flat fee.'' As an expression with
$n$ guests: $\mathbf{15n+40}$.

\textbf{Section I --- Order of Operations.}\\
34. $36\div(2+4^2)=36\div(2+16)=36\div18=\mathbf{2}$.\\
35. $5^2-3\times(8-2)=25-3\times6=25-18=\mathbf{7}$.\\
36. $(9-4)^2+3\times2=5^2+6=25+6=\mathbf{31}$.\\
37. $48\div(2^3-2)+5=48\div(8-2)+5=48\div6+5=8+5=\mathbf{13}$.

\textbf{Section J --- Evaluating Algebraic Expressions.}\\
38. $4(2)\cdot3\div6 = 24\div6 = \mathbf{4}$.\\
39. $(3\cdot2)^2\div(5-1) = 6^2\div4 = 36\div4 = \mathbf{9}$.\\
40. $2(4+3)-5 = 2(7)-5 = 14-5 = \mathbf{9}$.

\textbf{Section K --- Apply: Word Problems.}\\
41. Area $=\$960\div\$8\text{ per ft}^2 = 120$ ft$^2$. Side $=\sqrt{120}\approx10.95$, so the room is about $\mathbf{11}$ feet on a side.\\
42. Length $=\sqrt{6^2+(3\cdot6)^2}=\sqrt{36+324}=\sqrt{360}\approx18.9736\dots$, which rounds to $\mathbf{19.0}$ feet (nearest tenth).\\
43. Cost $=2(48)+3(27)+6h = 96+81+6h = \mathbf{177+6h}$. At $h=4$: $177+6(4)=177+24=\mathbf{\$201}$.

\end{document}
